\section{Conclusions and lessons learned}

\section{Limitations found}
The main drawbacks I have found during the development and research process during the whole work have been mainly hardware-related as the specific software requirements necessary to carry on the development are severely demanding due to the fact that the majority of the processes are run and rendered in the GPU. Although I have had access to a good graphical processing unit, the other components of the my working device have fallen behind. Taking for instance the minimum requirements of the used framework \textit{IsaacSim}, I found myself in the limits of the bare minimum, and some of my components such as the CPU were worse than the required. However, I have been able to carry on anyways although there have been a few moments where the I have really felt the bottleneck. When it comes those moments I have felt limited and slowed downed due to my equipment are the following:

\begin{itemize}
	\item I have needed to split some of the working components of the pipeline process inside \textit{IsaacSim} as the whole process would not fit in my GPU's VRAM. The simple fact of opening the framework already consumed close to a quarter of the graphic card's total capacity, and starting the server to listen and interact with the VLM calls alongside left little to no room for more processes in the background. The fact of working on the edge of the minimum hardware requirements also slowed the rendering of visual work and reduced the number of frames per second, nevertheless, this did not obstruct the development process.
	
	\item When it comes to making the most of the reinforcement learning training process, I have needed to tweak some of the predefined system configurations to maximize the performance and reduce any kind of lagging and overflow. Adapting the GPU capacity properties to maximize the overall functioning. If I were to work with better equipment, I would have liked to go for longer and deeper trainings, however, I am happy with the performance I obtained when tweaking the parameters I will have discussed in the corresponding section and other strategies such as parallelization.
	
	\item Not all limitations have been hardware related, as the official NVIDIA documentation websites were usually corrupted depending on the version. Leading to consulting external pages to solve the mismatching dependencies or installation processes for the frameworks used.
	
	\item After consulting the official documentation of the \textit{IsaacLab} repository, I found that even the default reinforcement learning training examples and demonstrations did not work as intended independently of the version. I was not the only one with the same issues as other peoples raised the same issues in different forums, so I inspired in the solutions given by the users with the same problems.
\end{itemize}

\section{Future work}